% sec_strip_10_hardy_aux_lattice.tex -- inactive archival fragment.
% Not part of the active article source graph rooted at main.tex.
% Deliberately inert unless prepared as a separate supplement.
\endinput

\subsection{Hardy splitting and the integer-orbit projection}

\begin{proposition}[Classical Hardy splitting of the Poisson strip model]%
\label{prop:poisson-hardy-splitting}
Let
\[
\mathcal H_K^+
:=
\left\{
f\in\mathcal H_K:\ \widehat f(\xi)=0 \text{ for a.e. }\xi<0
\right\},
\]
\[
\mathcal H_K^-
:=
\left\{
f\in\mathcal H_K:\ \widehat f(\xi)=0 \text{ for a.e. }\xi>0
\right\}.
\]
Then the classical positive/negative Hardy decomposition transported through
the Poisson image model gives
\[
\mathcal H_K=\mathcal H_K^+\oplus \mathcal H_K^-
\]
orthogonally.  Moreover, each $f_\pm\in\mathcal H_K^\pm$ extends
holomorphically to $\mathbb S$, and $\mathcal H_K^\pm$ is a
reproducing-kernel Hilbert space on~$\mathbb S$ with kernel
\begin{equation}\label{eq:hardy-splitting-kernels}
K_+(z,w)=\frac{1}{2(2-\ii(z-\overline w))},
\qquad
K_-(z,w)=\frac{1}{2(2+\ii(z-\overline w))}.
\end{equation}
For $a,b\in\RR$ one has
\begin{equation}\label{eq:hardy-splitting-real-kernel}
K(a,b)=K_+(a,b)+K_-(a,b).
\end{equation}
Thus the proposition is a background identification of the two Hardy
half-strip components of the already constructed Poisson image space; the
Poisson-specific data retained for the lattice theorem are the normalization
$K(a,b)=\pi P_2(a-b)$ and the explicit kernels in
\eqref{eq:hardy-splitting-kernels}.
\end{proposition}

\begin{proof}
This is the standard Paley--Wiener/Hardy support splitting for a strip,
written in the normalization of Proposition~\ref{prop:gram-space-poisson-image};
see the classical Hardy-space references in
Remark~\ref{rem:poisson-image-framework}.
The norm formula from Proposition~\ref{prop:gram-space-poisson-image}
shows that $\mathcal H_K^+$ and $\mathcal H_K^-$ are closed subspaces of
$\mathcal H_K$, and they are orthogonal because their Fourier supports are
disjoint.  The decomposition
\[
\widehat f
=\widehat f\,\mathbf 1_{[0,\infty)}
+\widehat f\,\mathbf 1_{(-\infty,0]}
\]
therefore yields the asserted orthogonal sum.

For $w\in\mathbb S$, define $K_\pm(\cdot,w)$ by
\[
\widehat{K_+(\cdot,w)}(\xi)
:=
\pi e^{-2\xi}e^{-\ii\overline w\,\xi}\mathbf 1_{[0,\infty)}(\xi),
\]
\[
\widehat{K_-(\cdot,w)}(\xi)
:=
\pi e^{2\xi}e^{-\ii\overline w\,\xi}\mathbf 1_{(-\infty,0]}(\xi).
\]
Since $|\Im w|<1$, these functions belong to $\mathcal H_K^\pm$.  If
$f_+\in\mathcal H_K^+$, then
\begin{align*}
\langle f_+,K_+(\cdot,w)\rangle_{\mathcal H_K}
&=
\frac{1}{2\pi^2}\int_0^\infty
e^{2\xi}\widehat f_+(\xi)\,
\overline{\pi e^{-2\xi}e^{-\ii\overline w\,\xi}}\dd\xi \\
&=
\frac{1}{2\pi}\int_0^\infty \widehat f_+(\xi)e^{\ii w\xi}\dd\xi
=f_+(w).
\end{align*}
The same argument gives the reproducing identity for $\mathcal H_K^-$.  By
Fourier inversion,
\[
K_+(z,w)
=\frac{1}{2\pi}\int_0^\infty
\pi e^{-2\xi}e^{\ii(z-\overline w)\xi}\dd\xi
=\frac{1}{2(2-\ii(z-\overline w))},
\]
\[
K_-(z,w)
=\frac{1}{2\pi}\int_{-\infty}^0
\pi e^{2\xi}e^{\ii(z-\overline w)\xi}\dd\xi
=\frac{1}{2(2+\ii(z-\overline w))},
\]
which proves \eqref{eq:hardy-splitting-kernels}.  Adding the two kernels
on the real axis yields \eqref{eq:hardy-splitting-real-kernel}.
\end{proof}

We next record the integer-orbit projection theorem used later.  This is the
only lattice projection needed for the observation-channel result.  We do not
introduce a variable spacing: throughout this subsection the lattice is
\(\ZZ\).  The calculation fixes the periodized Poisson weight, the symbol
\(\sigma\), and the orthogonal projection onto the closed span of the integer
kernel orbit.

\begin{lemma}[Integer projection for the Poisson strip kernel]%
\label{lem:auxiliary-uniform-lattice-projection}
Set
\[
I:=[-\pi,\pi],
\qquad
\Sigma_1(\theta):=\sum_{m\in\ZZ}e^{-2\abs{\theta+2\pi m}}
\qquad(\theta\in I).
\]
For
\[
K(a,b):=\frac{2}{4+(a-b)^2},
\qquad k(x):=K(x,0),
\]
one has
\[
\widehat k(\xi)=\pi e^{-2\abs{\xi}},
\]
and hence the integer Toeplitz symbol is
\begin{equation}\label{eq:integer-projection-poisson-symbol}
\sigma(\theta):=\sum_{n\in\ZZ}K(n,0)e^{-\ii n\theta}
=\pi\sum_{m\in\ZZ}e^{-2\abs{\theta+2\pi m}}
=\pi\frac{\cosh(2(\pi-\abs{\theta}))}{\sinh(2\pi)}
\qquad(\theta\in I).
\end{equation}
Equivalently,
\begin{equation}\label{eq:lattice-periodized-weight}
\Sigma_1(\theta)
=\frac{\cosh(2(\pi-\abs{\theta}))}{\sinh(2\pi)}
\qquad(\theta\in I).
\end{equation}
Define
\[
\mathcal H_K(\ZZ)
:=
\overline{\operatorname{span}}\{K(\cdot,n):n\in\ZZ\}
\subset\mathcal H_K
\]
and
\[
\mathcal N_{\ZZ}
:=
\{f\in\mathcal H_K:\ f(n)=0\ \text{for every }n\in\ZZ\}.
\]
For $f\in\mathcal H_K$ write
\[
q_f(\xi):=e^{2\abs{\xi}}\widehat f(\xi).
\]
Then the following hold.
\begin{enumerate}[label=(\roman*)]
\item
      One has $f\in\mathcal H_K(\ZZ)$ if and only if $q_f$ is
      $2\pi$-periodic and
      \[
      \int_I\Sigma_1(\theta)\abs{q_f(\theta)}^2\dd\theta<\infty.
      \]
      In that case,
      \begin{equation}\label{eq:lattice-subspace-norm}
      \|f\|_{\mathcal H_K}^2
      =\frac{1}{2\pi^2}\int_I\Sigma_1(\theta)\abs{q_f(\theta)}^2\dd\theta.
      \end{equation}
\item
      The orthogonal complement of~$\mathcal H_K(\ZZ)$ is exactly
      $\mathcal N_{\ZZ}$, and
      \begin{equation}\label{eq:lattice-alias-cancellation}
      f\in\mathcal N_{\ZZ}
      \iff
      \sum_{m\in\ZZ}\widehat f(\theta+2\pi m)=0
      \quad\text{for a.e. }\theta\in I.
      \end{equation}
      Consequently,
      \[
      \mathcal H_K=\mathcal H_K(\ZZ)\oplus \mathcal N_{\ZZ}
      \]
      orthogonally.
\item
      Let $\Pi_{\ZZ}:\mathcal H_K\to\mathcal H_K(\ZZ)$ denote the
      orthogonal projection.  Then, for almost every $\theta\in I$, the
      series
      \[
      \sum_{m\in\ZZ}\widehat f(\theta+2\pi m)
      \]
      converges absolutely and
      \begin{equation}\label{eq:lattice-alias-projection}
      q_{\Pi_{\ZZ}f}(\theta)
      =\frac{\displaystyle\sum_{m\in\ZZ}
      \widehat f(\theta+2\pi m)}
      {\Sigma_1(\theta)},
      \end{equation}
      equivalently
      \begin{equation}\label{eq:integer-orbit-projection-symbol}
      q_{\Pi_{\ZZ}f}(\theta)
      =
      \frac{\displaystyle\sum_{m\in\ZZ}\widehat f(\theta+2\pi m)}
      {\displaystyle\sum_{m\in\ZZ}e^{-2|\theta+2\pi m|}},
      \qquad -\pi\le\theta\le\pi,
      \end{equation}
      with $q_{\Pi_{\ZZ}f}$ extended $2\pi$-periodically to~$\RR$.
      Equivalently,
      \[
      \widehat{\Pi_{\ZZ}f}(\xi)=e^{-2\abs{\xi}}q_{\Pi_{\ZZ}f}(\xi)
      \qquad\text{for a.e. }\xi\in\RR.
      \]
      Moreover,
      \begin{equation}\label{eq:lattice-projection-norm}
      \|\Pi_{\ZZ}f\|_{\mathcal H_K}^2
      =\frac{1}{2\pi^2}\int_I
      \frac{\left|\sum_{m\in\ZZ}
      \widehat f(\theta+2\pi m)\right|^2}
      {\Sigma_1(\theta)}\,\dd\theta,
      \end{equation}
      and $\Pi_{\ZZ}f$ is the unique element of~$\mathcal H_K(\ZZ)$
      satisfying
      \[
      (\Pi_{\ZZ}f)(n)=f(n)
      \qquad(n\in\ZZ).
      \]
\item
      The integer samples determine exactly the projection $\Pi_{\ZZ}f$ and
      no finer ambient datum: for $f,g\in\mathcal H_K$,
      \begin{equation}\label{eq:uniform-lattice-quotient-rigidity}
      f(n)=g(n)\quad(n\in\ZZ)
      \qquad\Longleftrightarrow\qquad
      \Pi_{\ZZ}f=\Pi_{\ZZ}g .
      \end{equation}
      Equivalently,
      \[
      \mathcal H_K/\mathcal N_{\ZZ}\longrightarrow\mathcal H_K(\ZZ),
      \qquad
      f+\mathcal N_{\ZZ}\longmapsto\Pi_{\ZZ}f,
      \]
      is a well-defined isometric Hilbert-space isomorphism.  Thus full
      reconstruction from the integer orbit holds for a given
      $f\in\mathcal H_K$ if and only if $f\in\mathcal H_K(\ZZ)$.
\item
      Any reconstruction rule
      \[
      \mathcal R:\mathcal H_K\longrightarrow \mathcal H_K(\ZZ)
      \]
      whose value depends only on the integer samples and which is exact on
      $\mathcal H_K(\ZZ)$ must equal the integer projection:
      \[
      \mathcal Rf=\Pi_{\ZZ}f\qquad(f\in\mathcal H_K).
      \]
      Thus the integer orbit has a single possible ambient reconstruction:
      the orthogonal projection.
\end{enumerate}
\end{lemma}

\begin{proof}
The transform identity
\(\widehat k(\xi)=\pi e^{-2\abs{\xi}}\) is the standard Fourier transform
of \(2/(4+x^2)\) in the convention of this paper.  Poisson summation gives
\[
\sum_{n\in\ZZ}K(n,0)e^{-\ii n\theta}
=\sum_{m\in\ZZ}\widehat k(\theta+2\pi m)
=\pi\sum_{m\in\ZZ}e^{-2\abs{\theta+2\pi m}},
\]
which is the first equality in \eqref{eq:integer-projection-poisson-symbol}.
For $\theta\in I$,
\begin{align*}
\Sigma_1(\theta)
&=e^{-2\abs{\theta}}
+\sum_{m=1}^{\infty}e^{-2(\abs{\theta}+2\pi m)}
+\sum_{m=1}^{\infty}e^{-2(2\pi m-\abs{\theta})} \\
&=\frac{e^{-2\abs{\theta}}+e^{-4\pi}e^{2\abs{\theta}}}{1-e^{-4\pi}}
=\frac{\cosh(2(\pi-\abs{\theta}))}{\sinh(2\pi)},
\end{align*}
which gives both \eqref{eq:integer-projection-poisson-symbol} and
\eqref{eq:lattice-periodized-weight}.

If
\[
F=\sum_{n=-N}^{N}c_nK(\cdot,n),
\]
then
\[
\widehat F(\xi)=\pi e^{-2\abs{\xi}}\sum_{n=-N}^{N}c_ne^{-\ii n\xi},
\]
hence
\[
q_F(\xi)=\pi\sum_{n=-N}^{N}c_ne^{-\ii n\xi},
\]
which is $2\pi$-periodic.  Using the norm formula from
Proposition~\ref{prop:gram-space-poisson-image} and partitioning~$\RR$
into the translates $I+2\pi m$ gives
\[
\|F\|_{\mathcal H_K}^2
=\frac{1}{2\pi^2}\int_I
\sum_{m\in\ZZ}e^{-2\abs{\theta+2\pi m}}\abs{q_F(\theta)}^2\dd\theta
=\frac{1}{2\pi^2}\int_I\Sigma_1(\theta)\abs{q_F(\theta)}^2\dd\theta.
\]
Passing to the closure proves one direction of~(i).

Conversely, let $q$ be $2\pi$-periodic with
\[
\int_I\Sigma_1(\theta)\abs{q(\theta)}^2\dd\theta<\infty.
\]
Since $\Sigma_1$ is continuous and strictly positive on the compact
interval~$I$, the spaces
$L^2(I,\Sigma_1(\theta)\dd\theta)$ and $L^2(I,\dd\theta)$ coincide with
equivalent norms.  The trigonometric polynomials
\[
\sum_{n=-N}^{N}d_ne^{-\ii n\theta}
\]
are therefore dense in $L^2(I,\Sigma_1(\theta)\dd\theta)$.  Choose such a
sequence $p_N$ with $p_N\to \pi^{-1}q$ in that weighted norm, and define
$F_N\in\mathcal H_K(\ZZ)$ by
\[
F_N:=\sum_{n=-N}^{N}d_nK(\cdot,n).
\]
Then $q_{F_N}=\pi p_N$, so the preceding norm identity gives
\[
\|F_N-F_M\|_{\mathcal H_K}^2
=\frac{1}{2\pi^2}\int_I\Sigma_1(\theta)
\abs{q_{F_N}(\theta)-q_{F_M}(\theta)}^2\dd\theta.
\]
Hence $(F_N)$ is Cauchy in~$\mathcal H_K$ and converges to some
$F\in\mathcal H_K(\ZZ)$ with $q_F=q$.  This proves~(i).

For every $f\in\mathcal H_K$ and almost every $\theta\in I$,
Cauchy--Schwarz gives
\[
\sum_{m\in\ZZ}\abs{\widehat f(\theta+2\pi m)}
\le
\Sigma_1(\theta)^{1/2}
\left(
\sum_{m\in\ZZ}e^{2\abs{\theta+2\pi m}}
\abs{\widehat f(\theta+2\pi m)}^2
\right)^{1/2},
\]
so the alias sum is absolutely convergent for almost every~$\theta$.

If $g\in\mathcal H_K(\ZZ)$, then by~(i) its profile $q_g$ is
$2\pi$-periodic.  Using the Fourier norm formula again,
\begin{align*}
\langle f,g\rangle_{\mathcal H_K}
&=
\frac{1}{2\pi^2}\int_I
\sum_{m\in\ZZ}e^{-2\abs{\theta+2\pi m}}q_f(\theta+2\pi m)\,
\overline{q_g(\theta)}\dd\theta \\
&=
\frac{1}{2\pi^2}\int_I
\left(
\sum_{m\in\ZZ}\widehat f(\theta+2\pi m)
\right)\overline{q_g(\theta)}\dd\theta.
\end{align*}
Since the periodic profiles $q_g$ form a dense subspace of
$L^2(I,\Sigma_1(\theta)\dd\theta)$ and $\Sigma_1$ is bounded above and below
on~$I$, the same set is dense in $L^2(I,\dd\theta)$.  Therefore
$f\perp\mathcal H_K(\ZZ)$ if and only if the alias sum in
\eqref{eq:lattice-alias-cancellation} vanishes almost everywhere on~$I$.

By the reproducing property,
\[
f(n)=\langle f,K(\cdot,n)\rangle_{\mathcal H_K}
\qquad(n\in\ZZ),
\]
and the closed linear span of the kernel sections $K(\cdot,n)$ is
exactly $\mathcal H_K(\ZZ)$.  Consequently,
\[
f\in\mathcal N_{\ZZ}
\iff
f\perp\mathcal H_K(\ZZ).
\]
This proves~(ii).

Let $g\in\mathcal H_K(\ZZ)$ and write $q_g$ for its periodic profile.
Then
\[
\|f-g\|_{\mathcal H_K}^2
=\frac{1}{2\pi^2}\int_I
\sum_{m\in\ZZ}e^{-2\abs{\theta+2\pi m}}
\abs{q_f(\theta+2\pi m)-q_g(\theta)}^2\dd\theta.
\]
For each $\theta\in I$, the integrand is minimized when $q_g(\theta)$ is the
weighted average of the fiber values $q_f(\theta+2\pi m)$ with weights
$e^{-2\abs{\theta+2\pi m}}$.  Since
\[
e^{-2\abs{\theta+2\pi m}}q_f(\theta+2\pi m)
=\widehat f(\theta+2\pi m),
\]
the minimizing profile is exactly the one in
\eqref{eq:lattice-alias-projection}.  This proves the projection formula.
Since \(\Sigma_1(\theta)=\sum_m e^{-2\abs{\theta+2\pi m}}\), this is also
\eqref{eq:integer-orbit-projection-symbol}.
Substituting the minimizing value into the norm identity gives
\eqref{eq:lattice-projection-norm}.  Finally,
\[
f-\Pi_{\ZZ}f\in\mathcal N_{\ZZ}
\]
by~(ii), so $(\Pi_{\ZZ}f)(n)=f(n)$ for every $n\in\ZZ$.  If
$G\in\mathcal H_K(\ZZ)$ satisfies $G(n)=f(n)$ for all $n$, then
$G-\Pi_{\ZZ}f\in\mathcal H_K(\ZZ)\cap\mathcal N_{\ZZ}=\{0\}$, hence
$G=\Pi_{\ZZ}f$.  This proves~(iii).

It remains to identify the quotient determined by the samples.  The forward
implication in \eqref{eq:uniform-lattice-quotient-rigidity} follows because
equal integer samples give $f-g\in\mathcal N_{\ZZ}$, and by~(ii) this is the
orthogonal complement of $\mathcal H_K(\ZZ)$, so the two orthogonal
projections coincide.  The reverse implication follows from~(iii), since
$\Pi_{\ZZ}f$ and $\Pi_{\ZZ}g$ interpolate respectively the samples of
$f$ and~$g$.
The displayed quotient map is well defined because adding an element of
$\mathcal N_{\ZZ}=\mathcal H_K(\ZZ)^\perp$ does not change the projection.  It is
surjective by applying it to $F+\mathcal N_{\ZZ}$ with $F\in\mathcal H_K(\ZZ)$,
injective by \eqref{eq:uniform-lattice-quotient-rigidity}, and isometric by
the orthogonal decomposition
\[
\|f+\mathcal N_{\ZZ}\|_{\mathcal H_K/\mathcal N_{\ZZ}}
=\inf_{u\in\mathcal N_{\ZZ}}
\|\Pi_{\ZZ}f+(f-\Pi_{\ZZ}f)+u\|_{\mathcal H_K}
=\|\Pi_{\ZZ}f\|_{\mathcal H_K}.
\]
Finally, full reconstruction means $f$ equals the unique lattice-subspace
interpolant $\Pi_{\ZZ}f$, and this is equivalent to
$f\in\mathcal H_K(\ZZ)$.  This proves~(iv).

For~(v), if $\mathcal R$ depends only on the integer samples, then $f$ and
$\Pi_{\ZZ}f$ have the same samples by~(iii), so
\[
\mathcal Rf=\mathcal R(\Pi_{\ZZ}f).
\]
Since $\Pi_{\ZZ}f\in\mathcal H_K(\ZZ)$ and $\mathcal R$ is exact on that
subspace, the right-hand side is $\Pi_{\ZZ}f$.  Hence
$\mathcal R=\Pi_{\ZZ}$.  The final sentence records precisely the
integer-orbit input used by the later observation theorem.
This proves~(v).
\end{proof}

